\documentclass[11pt]{article}

\usepackage[margin=1in]{geometry}
\usepackage{amsmath,amssymb}
\usepackage{graphicx}
\usepackage{booktabs}
\usepackage{hyperref}
\usepackage{natbib}
\usepackage{authblk}
\usepackage{caption}

\title{Silent Attractors: Adversarial Task Design for Detecting\\
Latent Reasoning Failures in Frontier Language Models}

\author[1]{Polex Ouma Otieno}
\affil[1]{Independent Researcher / Smartlinks Ltd}

\date{\today}

\begin{document}
\maketitle

\begin{abstract}
Standard benchmarks increasingly saturate against frontier language
models, yet models continue to fail in production on tasks that require
one non-obvious reasoning step diverging from a trained professional
reflex. We introduce a task design methodology for constructing
\emph{adversarial evaluation tasks} that isolate this failure mode,
which we term a \textbf{silent attractor}: a plausible, confidently
stated answer produced by a common heuristic or memorized pattern that
is not supported by the specific facts of the problem. Each task is
validated against two references before inclusion: an oracle solution
that must pass, and a naive/reflex solution constructed from the most
common professional shortcut, which must fail. We describe a nine-family
taxonomy of failure mechanisms observed across accepted tasks spanning
financial accounting, pension cost restatement, and distributed systems
domains, and report pass rates of current frontier models against a
pilot task set. We release our task construction methodology, a
reproducible evaluation harness, and the pilot task set to support
further adversarial benchmark construction.
\end{abstract}

\section{Introduction}

Large language models are increasingly deployed in domains where a
wrong-but-confident answer carries real financial or operational cost:
accounting close processes, engineering sign-off, and regulatory
reporting among them. Aggregate benchmark accuracy, however, is a poor
proxy for reliability in these settings, because a model can score well
on average while reliably failing on the specific subset of problems
where a memorized professional reflex diverges from the fact pattern at
hand.

This paper describes a task construction methodology, developed and
iterated across an adversarial benchmark contribution platform, for
deliberately isolating this failure mode. We call it a \textbf{silent
attractor}: the model converges on an answer that would be correct
under the common case, but the specific problem instance has been
constructed so that one fact --- present in the prompt, never
flagged as unusual --- changes the correct derivation path. The model
is not being tricked with irrelevant information or adversarial
phrasing; every fact given is necessary and sufficient to solve the
problem correctly. The failure is purely a reasoning failure: applying
the high-frequency pattern instead of deriving the answer from the
specific facts given.

\textbf{Contributions.} (1) A task construction and validation
methodology --- oracle-passes / naive-reflex-fails --- for building
silent-attractor tasks. (2) A nine-family taxonomy of the reasoning
mechanisms that produce this failure mode. (3) A reproducible evaluation
harness (task loader, multi-provider model runner, automatic verifier,
and statistical reporting with Wilson score confidence intervals) release
alongside a pilot task set. (4) Pilot results across current frontier
models on financial accounting tasks constructed with this methodology.

\section{Related Work}

Adversarial and contamination-resistant benchmark construction has grown
alongside concerns that static benchmarks saturate and leak into
training data. Work on LLM-judge reliability has shown that automated
evaluators themselves exhibit systematic biases, including
overconfidence and reduced sensitivity to expert-identified weaknesses
compared to human raters, motivating evaluation designs where the
correct answer is verifiable by direct computation rather than judged
by another model. Separately, benchmark work on agentic and
multi-turn settings has shown that models are sensitive to
distributional shift away from the common case, including under peer
pressure from unreliable collaborators. Our work differs in scope: we
target single-turn, professionally grounded quantitative tasks in
which the shift away from the common case is entirely contained in the
stated facts, with no adversarial framing, multi-agent interaction, or
distributional shift required.

\section{Task Construction Methodology}

Each task begins from a real professional workflow (e.g., a pension
cost restatement, a foreign currency consolidation) and is built
through the following stages.

\begin{enumerate}
  \item \textbf{Identify the reflex.} Identify the single most common
    shortcut, heuristic, or memorized formula a trained professional
    would reach for given the surface features of the problem.
  \item \textbf{Identify the crux.} Identify one fact, derivable
    without any information beyond what is given, that changes the
    correct derivation path away from that reflex.
  \item \textbf{Construct the oracle.} Write a reference solution that
    correctly incorporates the crux.
  \item \textbf{Construct the naive baseline.} Write a second solution
    that applies the identified reflex and ignores the crux.
  \item \textbf{Validate discrimination.} Confirm the oracle answer and
    naive answer differ by a margin well outside the task's stated
    tolerance, so the task can actually discriminate between correct
    and merely-plausible reasoning.
  \item \textbf{Author the instruction.} Write the task prompt so that
    every fact needed for the oracle solution is present, unflagged,
    and stated plainly, with no hint that a shortcut would be
    incorrect.
  \item \textbf{Assemble source artifacts.} Where the underlying
    workflow naturally spans multiple documents (a data extract, a
    spreadsheet, a policy memo or plan document), construct each as a
    separate artifact rather than a single flattened prompt, so the
    task requires cross-referencing rather than single-document
    reading comprehension.
\end{enumerate}

This mirrors a validated production workflow: every task in the pilot
set was built and reviewed through a formal proposal and multi-stage
technical review pipeline before inclusion, including an automated
check that the oracle solution passes and the naive solution fails.

\subsection{Failure mechanism taxonomy}

Across accepted tasks, we identify nine recurring mechanisms by which a
silent attractor is constructed. We name each mechanism descriptively
rather than by an arbitrary label, since the name itself is part of
the contribution:

\begin{description}
  \item[Reflex override] A trained professional reflex produces the
    wrong answer; correct reasoning requires one non-obvious
    derivation step the data supports but never announces.
  \item[Heuristic substitution] A shortcut resembles the correct
    method closely enough to look right while silently diverging from
    it.
  \item[Distractor injection] A distractor resource or datum invites
    misuse if not filtered out.
  \item[Evidence synthesis gap] The answer must be inferred from
    partial evidence rather than read off directly.
  \item[Assumption violation] An assumption that normally holds is
    violated in this specific instance.
  \item[Weak-signal fusion] The failure only becomes visible once
    multiple individually weak signals are combined.
  \item[Exhaustive coverage failure] Many sub-parts must all be
    correct simultaneously for the task to be scored as passed.
  \item[Temporal coupling] Order-of-operations or path dependence
    changes the correct answer.
  \item[Staleness error] A stale or superseded value is used where the
    value as of a specific date is required.
\end{description}

The two pilot tasks released with this paper (Section~\ref{sec:pilot})
instantiate temporal coupling and assumption violation respectively.

\section{Evaluation Harness}

We release a lightweight, provider-agnostic evaluation harness
alongside this paper (see the accompanying code repository):

\begin{itemize}
  \item \textbf{Task loader} parses a task directory (\texttt{task.toml}
    + \texttt{instruction.md}) into a uniform schema, so task exports
    from the source platform require minimal reshaping.
  \item \textbf{Model runner} sends each task's instruction to one or
    more configured model providers and stores the raw completion.
    Anthropic and OpenAI clients are provided; adding a provider
    requires implementing a single \texttt{complete(prompt)} method.
  \item \textbf{Verifier} extracts a numeric final answer from the raw
    completion and scores it against the oracle answer within a
    per-task tolerance.
  \item \textbf{Failure tagger} attaches each scored result to its
    task's failure family, enabling aggregation by mechanism rather
    than only by task or model.
  \item \textbf{Statistical reporting} computes pass rates with Wilson
    score confidence intervals, appropriate for the small-sample
    binomial proportions typical of a hand-constructed adversarial task
    set, and renders bar-chart figures directly from scored results.
\end{itemize}

\section{Pilot Results}
\label{sec:pilot}

We report results on a two-task pilot set spanning ASC 715 (pension
accounting) and ASC 830 (foreign currency consolidation), instantiating
the temporal coupling and assumption violation mechanisms respectively.
Both tasks were validated per Section~3: the oracle solution passes
within a 1\% tolerance and the naive reflex solution fails by a wide
margin (24.3\% and 20.0\% relative error respectively). Each task is
accompanied by multiple source artifacts (a SQLite data extract, and
either a spreadsheet or a short policy memo PDF) mirroring the
multi-source structure of the source platform's tasks, fictionalized to
avoid disclosing any live evaluation content.

\begin{table}[h]
\centering
\caption{Pilot task summary.}
\label{tab:pilot-tasks}
\begin{tabular}{lllr}
\toprule
Task & Domain & Failure mechanism & Naive/oracle gap \\
\midrule
demo\_pension\_restatement & ASC 715 & Temporal coupling & 24.3\% \\
demo\_fx\_consolidation & ASC 830 & Assumption violation & 20.0\% \\
\bottomrule
\end{tabular}
\end{table}

\begin{figure}[h]
\centering
\includegraphics[width=0.75\textwidth]{../figures/pass_rate_by_family.png}
\caption{Pass rate by failure family (pilot run; replace with results
from a full model sweep before submission).}
\label{fig:pass-rate-family}
\end{figure}

\textit{Note: the figures and tables in this draft are generated from a
harness smoke test using a deterministic offline mock client, to
confirm the pipeline runs end to end. Before submission, re-run
\texttt{src/run\_pipeline.py} with \texttt{--provider anthropic} and
\texttt{--provider openai} (and any additional providers) against the
full task set to produce real results.}

\section{Discussion and Limitations}

The pilot set reported here is intentionally small; the central
contribution of this paper is the task construction methodology and
harness, not a large-scale leaderboard. A full submission should expand
the task set substantially (target: 30--50 tasks spanning multiple
failure families and domains) and report results across a representative
set of current frontier models. We also note that automatic numeric
answer extraction is a simplification; tasks requiring structured or
free-text answers would need a more general verifier, and any use of a
model-based judge for non-numeric answers should be cross-checked
against known LLM-judge reliability issues.

The current harness evaluates models via single-turn text completion:
each task's source artifacts (data extracts, spreadsheets, policy
documents) are flattened into the instruction text rather than provided
as files a tool-using agent must open and parse. This keeps the harness
simple and provider-agnostic, but it understates the difficulty of the
source platform's actual tasks, which are agentic and require an
agent to independently locate, parse, and cross-reference multiple raw
files. Extending the harness to an agentic, tool-using evaluation loop
(e.g. giving the model file-system and code-execution access rather
than a single flattened prompt) is the highest-value next step for
making pilot results comparable to the source platform's own difficulty
gate.

\section{Conclusion}

We present a methodology and reproducible harness for constructing
adversarial evaluation tasks that isolate silent-attractor reasoning
failures in frontier language models, together with a nine-family
taxonomy of the mechanisms by which such failures are constructed. We
release a pilot task set and invite extension to additional domains and
failure families.

\bibliographystyle{plainnat}
\bibliography{references}

\end{document}
